\documentclass[10pt,twoside]{book}
\usepackage[utf8]{inputenc}
\usepackage[T1]{fontenc}
\usepackage{geometry}
\geometry{left=1.2cm,right=1.2cm,top=1.5cm,bottom=1.5cm}
\usepackage{amsmath,amssymb}
\usepackage{graphicx}
\usepackage{booktabs}
\usepackage{tabularx}
\usepackage{array}
\usepackage{xcolor}
\usepackage{titlesec}
\usepackage{setspace}
\usepackage{fancyhdr}
\usepackage{mdframed}
\usepackage{needspace}
\usepackage{enumitem}
\usepackage{hyperref}
\hypersetup{colorlinks=true,linkcolor=blue,urlcolor=blue}

\newcommand{\tag}[1]{\textsc{#1}}
\newcommand{\clk}[1]{\textit{[#1]}}
\newcommand{\stringitem}[1]{\textit{#1}}
\newcommand{\dusana}[1]{\begin{quote}\itshape #1 \\ --- Dusana of the Raven Road\end{quote}}

\titleformat{\chapter}{\Huge\bfseries}{\thechapter}{20pt}{}
\titleformat{\section}{\Large\bfseries}{\thesection}{15pt}{}
\setlength{\parindent}{0pt}
\setlength{\parskip}{0.8ex}

\newmdenv[linecolor=black,backgroundcolor=gray!5,innertopmargin=8pt,innerbottommargin=8pt,innerrightmargin=8pt,innerleftmargin=8pt]{infobox}

\begin{document}

\title{\Huge\textbf{The Salt Cycle}\\ \large Tales of the Threshold, from the Southern Roads}
\author{Told by the fires of a Sidhi lighthouse\\ \small Set down by Dusana of the Raven Road}
\date{Under the lamp, bread and salt --- and a promise that must be returned in stories}
\maketitle

\tableofcontents

\chapter*{Preface: The Storyteller's Road}
\addcontentsline{toc}{chapter}{Preface: The Storyteller's Road}

A cycle of stories is a cycle of promises. Some promises are written in ink. Some are written in salt. And some --- the oldest kind --- are written in nothing but the space between a teller's words.

I left the baronet's forest with a new wheel and an old promise. The Pereshi storyteller Kaveh had given me a tale in exchange for a cup of tea, years ago in the port of Zakov. I had promised to pass it on. I had not. The promise sat in my chest like a swallowed stone.

The road south is not kind to those who forget.

I crossed the Dolmis on a Sidhi trader's dhow, the captain a woman with one eye and a brass hook for a hand. She asked no fare --- only a story. I told her the one about the Fox's Wedding. She laughed and said, "The foxes in our waters are different. They wear shells, not silk. But a gift is a gift. You may have passage."

I walked through the Pereshi highlands, where the mountains are older than any empire and the djinn leave their footprints in the scree. I followed the spice roads into the southern lands --- not one land, but many. The weigh‑houses of \textbf{Oshiira}, where grain is measured against truth. The black marble gates of \textbf{Ashaan}, where the Veiled Sisters count names, not coin. The rolling paths of the \textbf{Ngomebe} walking cities, where the stone itself sings of broken vows. The highland courts of \textbf{Taharka}, where a horse's speed can settle a kingdom's fate. The honeyed groves of \textbf{Sekogo}, where the Green Council hears the spirits of trees as witnesses. The \textbf{Crimson Valley}, where the Lethai‑ar weave their promises into silk and cut the thread when a vow is forgotten.

And I listened to the small peoples --- the \textbf{Aelinnel} and the \textbf{Aelaerem} --- who live in the cracks between these powers, borrowing each other's customs and arguing over whether to count the ninth step or pour the ninth cup.

I have written nine tales. They are not the tale I promised. They are the weight that grew while I waited. This is the third cycle. The first was the Hearth Cycle, told in the north. The second was the Wandering Cycle, told in the east. This is the Salt Cycle, told on the southern roads --- where the sea is warm, the mountains are sharp, and the promise is still waiting to be returned.

The ink is squid. The lamp is borrowed. The road is south.

\dusana{At the crossroads of three empires, with a full waterskin and an empty purse. The ninth cup remains un-poured.}

\chapter{The First Night: The Lighthouse of the One-Eyed Captain}

\section*{Told by}
Captain Rana of the \textit{Eel-Skiff}, the same woman who had carried me across the Dolmis. She told this tale at anchor, with her brass hook resting on the rail.

\section*{The Tale}
The Sidhi say that every lighthouse was built over a drowned sailor's grave. Not the famous sailors --- the forgotten ones. The ones who washed ashore with no name and no country.

A lighthouse keeper --- call her Amara --- tended a beacon on the Sidhi coast. She was old, blind in one eye, and she spoke to the light as if it were her daughter.

"Why do you talk to the lamp?" a visitor asked.

"Because the lamp talks back," Amara said. "It tells me which ships are honest and which are carrying stolen salt."

The visitor laughed. Amara said, "Wait until midnight."

At midnight, a ship appeared on the horizon. Its sails were furled, but it moved against the wind. Amara turned the lamp a quarter degree south.

"That ship is a slaver," she said. "The light will lead it onto the reef."

The ship struck. The crew drowned. The slaves swam to shore.

The visitor was horrified. "You killed them!"

"I killed the slavers," Amara said. "The lamp chose. I only turned it."

She was arrested by the Admiralty. The magistrate asked, "How do you know the lamp speaks?"

Amara touched her blind eye. "Because I lost this eye to a lamp that would not speak for me. It saw a truth I did not want to see. So I cut the truth out. But the lamp --- the lamp forgave me."

The magistrate freed her. The lamp still burns on that headland. And every Sidhi sailor knows: if the light shifts a quarter degree south, you reef your sails and pray you are not carrying stolen salt.

\section*{The Witch's Closing}
"A lamp that speaks is a witness. A witness that judges is a court. Do not ask for justice unless you are prepared to be judged."

\section*{What Lurks in the Shadow of the Tale}
\begin{infobox}
\textbf{The Speaking Beacon (Sidhi variant)} --- A lighthouse whose lamp is alive. It can see the cargo of any ship on the horizon and will shift its beam to guide the worthy and wreck the guilty.

\textbf{Tags / Mechanics:} \tag{LOCATION} \tag{TRUTH} \tag{JUDGMENT}

\begin{itemize}
\item A character who tends the lamp for a night may ask one question about a ship visible from the shore. The lamp answers truthfully (GM discretion, but the answer must be useful).
\item If the character asks the lamp to wreck a ship, they must succeed on a \textit{Spirit + Lore} test (DV 4) to interpret the lamp's judgment. On a miss, the lamp wrecks a different ship --- possibly one the character intended to save.
\item \textbf{Clock: The Drowned's Reckoning [6]} -- ticks each time the lamp is used to judge. When full, a ghost ship appears, crewed by those the lamp has wrecked. The ghost ship demands an accounting.
\end{itemize}

\textbf{Adventure Hook:} A Sidhi lighthouse keeper has gone blind in both eyes. The lamp has gone silent. Ships are wrecking on a clear night. The party must discover what the lamp last saw --- and why it has stopped speaking.
\end{infobox}

\chapter{The Second Night: The Carpet That Remembers}

\section*{Told by}
Soraya the Weaver, a Pereshi woman whose family had woven carpets for seven generations. She spoke while her loom clicked, never looking up.

\section*{The Tale}
The Pereshi say that a carpet is not a thing. It is a covenant between the weaver and the land. The wool is a promise. The dye is a memory. The pattern is a map of unsettled weights.

A prince --- call him Darius --- commissioned a carpet from the greatest weaver in the mountains. He wanted a carpet that would remember his victories. The weaver agreed.

For a year, she wove. She used wool from a lamb born under a blood moon. She used dye from a flower that blooms only once a century. She wove the prince's battles, his hunting trips, his coronation.

When she finished, the prince was furious. "Where is my face? I am not in the carpet."

The weaver said, "You are in every thread. But the carpet does not show faces. It shows the weight of what was taken without return."

The prince looked closer. The battles were still there, but they were not his victories. They were the victories of the men who had died for him. Their faces were woven into the border --- small, barely visible, but present.

"You have woven my shame," the prince said.

"I have woven the truth," the weaver said. "A carpet that lies is a carpet that unravels."

The prince kept the carpet. He did not hang it. He stored it in a chest. But every night, he dreamed of the faces in the border. And every morning, he gave a coin to a widow.

When he died, they unrolled the carpet one last time. The faces were gone. In their place was a single word, woven in gold thread: \textit{Returned}.

\section*{The Witch's Closing}
"A carpet is a covenant. Walk on it, and you agree to its terms. Read the border before you step."

\section*{What Lurks in the Shadow of the Tale}
\begin{infobox}
\textbf{The Carpet of Weights} --- A Pereshi carpet that records the unacknowledged weight a person carries. It cannot lie. If the weight is released, the images fade. If the weight grows, the carpet expands.

\textbf{Tags / Mechanics:} \tag{ARTIFACT} \tag{TRUTH} \tag{CURSE}

\begin{itemize}
\item A character who sleeps on the carpet for a night sees a vision of one unsettled promise (material, emotional, or magical). The GM reveals a narrative hook.
\item The character may choose to release that promise. If they do, they gain +1 die to all social rolls involving the community for a week. If they refuse, the carpet's pattern shifts to show the character's face in the border, and they suffer --1 die to all social rolls until the promise is settled.
\item \textbf{Clock: The Unraveling [8]} -- ticks each time a promise is ignored. When full, the carpet tears itself apart, and the character's name is erased from their family's records.
\end{itemize}

\textbf{Adventure Hook:} A Pereshi merchant's carpet has grown to fill an entire room. He does not know whose unsettled weights are woven into it. The party must discover the faces in the border --- and find a way to release them before the carpet swallows the house.
\end{infobox}

\chapter{The Third Night: The Djinn's Three Questions}

\section*{Told by}
A nameless Hyaar goatherd who lived in a cave and had not spoken to a human in ten years. He spoke because I offered him salt.

\section*{The Tale}
The Hyaar say that the djinn live in the spaces between mountains --- not in the peaks, not in the valleys, but in the shadow of a rock at noon.

A young woman --- call her Zara --- got lost in the high passes. She was not afraid. She had grown up in the mountains. But the fog came down thick and gray, and she could not see her hand in front of her face.

A voice spoke from the mist. "You are in my house. You carry a question you have not asked."

Zara knew the rule. Djinn ask three questions. You must answer truthfully, or they keep you. She waited.

"First question," said the djinn. "What is the name of the wind that carries the snow?"

"The north wind," Zara said.

"Wrong," said the djinn. "The snow carries itself. The wind only watches. One answer left."

Zara thought. She had been tricked. The djinn did not want facts. It wanted truths that hurt.

"Second question," said the djinn. "Who died so that you could be born?"

Zara thought of her mother, who had died in childbirth. "My mother," she said.

"Correct," said the djinn. "Now the third. What are you willing to leave behind to find your way home?"

Zara thought. She thought of her father, who was old and alone. She thought of her dowry, which was three goats and a copper pot. She thought of her name, which she had never liked.

"My fear," she said.

The fog parted. The path was clear. The djinn was gone.

She walked home. She was not afraid of the mountains anymore. But sometimes, on the high passes, she felt the djinn watching --- not with malice, but with curiosity. It had never heard that answer before.

\section*{The Witch's Closing}
"The djinn do not want your gold. They want your certainty. Give them doubt, and they will let you pass."

\section*{What Lurks in the Shadow of the Tale}
\begin{infobox}
\textbf{The Djinn of the High Pass} --- A spirit of the mountains that appears in fog. It asks three questions. The first is a trick. The second is a wound. The third is a choice. Answer well, and it guides you. Answer poorly, and it keeps you.

\textbf{Tags / Mechanics:} \tag{SPIRIT} \tag{TRUTH} \tag{DREAD}

\begin{itemize}
\item The djinn appears when a character is lost and the fog is thick. The GM asks three questions, which the player must answer in character.
\item The first question must have a wrong answer. The player chooses whether to give the wrong answer (the djinn gains 1 SB) or the right answer (the character gains +1 die on their next roll).
\item The second question must touch on a character's backstory (e.g., "Who died so that you could be born?"). The player must answer truthfully or the djinn keeps them (the character is lost for 1d4 hours and marks 1 Fatigue).
\item The third question offers a choice. The player names something they are willing to lose. The djinn may or may not take it. On a successful \textit{Spirit + Resolve} test (DV 5), the djinn accepts the offering and guides the character home. On a failure, the djinn takes something else --- a memory, a skill point, or a relationship.
\end{itemize}

\textbf{Adventure Hook:} A Hyaar village is disappearing. Every week, someone gets lost in the fog. They return --- but they are different. They have answered the djinn's questions poorly. The party must find the djinn and negotiate for the village's memories.
\end{infobox}

\chapter{The Fourth Night: The Dwarf's Bridge}

\section*{Told by}
An Aeler tunnel-wright named Vorn, who had built bridges under three mountains. He smelled of stone dust and old beer.

\section*{The Tale}
The Hyaar Peninsula is split by a river that has no name. The dwarves call it the Vein. The Pereshi call it the Weight. The Sidhi call it the Salt Line.

A bridge was built across the Salt Line by an Aeler engineer --- call her Helga. She built it to last a thousand years. She carved her name into the keystone.

The bridge lasted a hundred. Then a flood washed away the approaches. The dwarves came to rebuild.

"The agreement says we built it to last a thousand years," the locals said.

"The agreement says nothing about floods," the dwarves said. "The flood is an act of the river. The river is not a signatory."

A Pereshi judge was called. She listened to both sides. Then she asked, "Where is the keystone?"

They brought the keystone. Helga's name was still carved into it. But underneath, in smaller letters, was a new name: \textit{The River}.

"The engineer signed on behalf of the river," the judge said. "The river did not consent. Therefore, the agreement is void. Rebuild the bridge. This time, ask the river for permission."

They did not know how to ask a river for permission. So they sent a Sidhi diver down to the riverbed. He came back with a single stone, smooth and black.

"The river says yes," he said. "It also says the toll is one story per crossing."

The dwarves rebuilt the bridge. They left a niche in the arch for a storyteller. And every traveler who crosses must tell a story to the river --- or pay double.

\section*{The Witch's Closing}
"A bridge is a conversation. Speak to the river before you build. It has a longer memory than your agreement."

\section*{What Lurks in the Shadow of the Tale}
\begin{infobox}
\textbf{The Salt Line Bridge} --- A dwarven bridge over a river that is a living presence. It demands a toll: one story per crossing. If the story is good, the bridge lets you pass. If the story is poor, the bridge grows a new crack --- and the next crossing costs double.

\textbf{Tags / Mechanics:} \tag{LOCATION} \tag{THRESHOLD} \tag{TRUTH}

\begin{itemize}
\item To cross, a character must tell a story (real or invented) to the river. The GM judges the story's quality. On a good story (genuine, clever, or moving), the bridge grants +1 die to the next action on the far side. On a poor story (repetitive, boring, or false), the bridge's toll for the return crossing is doubled (2 SB or 2 Fatigue).
\item A character may choose to pay double without telling a story. This marks 1 SB and 1 Fatigue.
\item \textbf{Clock: The River's Reckoning [6]} -- ticks each time a poor story is told. When full, the bridge collapses. The river demands a life --- or a new story good enough to raise the bridge again.
\end{itemize}

\textbf{Adventure Hook:} The Salt Line Bridge has collapsed. The dwarves refuse to rebuild without a new understanding. The river demands a story that has never been told. The party must find such a story --- or forge a new agreement with the river's spirit.
\end{infobox}

\chapter{The Fifth Night: The Sidhi Freedman's Seal}

\section*{Told by}
An old Sidhi dock-worker in Port Shed, who showed me a scar on his wrist where a slave collar had been burned off.

\section*{The Tale}
The Sidhi were slaves once. Not all of them, not forever --- but long enough that every family has a story of the collar.

A young Sidhi man --- call him Elias --- bought his freedom from an Ashaani master. He paid in silver and promises. The master gave him a seal --- a brass disk stamped with the words \textit{Freed by law}.

"This seal is your proof," the master said. "Lose it, and you are a slave again."

Elias wore the seal on a cord around his neck. He became a merchant. He prospered. He bought a ship. He hired other freedmen.

One night, a storm sank his ship. Elias swam to shore. The seal was gone.

He searched the beach for three days. Nothing.

He went to the Ashaani master. "I need a new seal."

The master laughed. "You are not a freedman without the seal. You are a slave who ran away. The law says I can take you back."

Elias did not run. He knelt. "Then take me. But let me say goodbye to my wife."

The master agreed. Elias went home. He told his wife, "Gather the crew. Tell them I am a slave again."

His wife said, "You are not a slave. You are a Sidhi. The law does not define you. The sea does."

She cut a piece of rope from the anchor line. She tied it around his neck. "This is your new seal. It is not brass. It is memory. Every knot is a freedman's name. When you have enough knots, the law will follow."

Elias returned to the Ashaani master. "I have no seal. But I have a rope with nine knots. Each knot is a witness. The law cannot enslave a man with nine witnesses."

The master looked at the rope. He looked at Elias. He said, "The law can do anything. But I am tired of this argument. Go. Keep your rope."

Elias became a legend. The Sidhi still tie ropes with nine knots --- one for each of the original witnesses --- and they wear them under their shirts. The law does not ask to see them. The law is afraid of what the knots remember.

\section*{The Witch's Closing}
"A seal is a promise. A memory is a witness. The law fears the witness more than the promise."

\section*{What Lurks in the Shadow of the Tale}
\begin{infobox}
\textbf{The Nine-Knot Rope} --- A length of rope with nine knots, each representing a witness to a promise. It is not a legal document, but it carries the weight of memory. Those who wear it cannot be easily enslaved --- the law hesitates to challenge the knots.

\textbf{Tags / Mechanics:} \tag{ARTIFACT} \tag{TRUTH} \tag{PROTECTION}

\begin{itemize}
\item A character who wears the Nine-Knot Rope gains +1 die to resist any attempt to enslave, imprison, or legally bind them against their will. The knots remind the law of witnesses.
\item Once per session, the character may untie one knot to "spend" a witness. The character gains +2 dice on a single social roll involving justice or freedom. The knot is gone; the witness is used.
\item \textbf{Clock: The Witness's Silence [4]} -- ticks each time a knot is used. When full, the character's name is removed from the legal records of the town where the knot was tied. No one remembers them there. They are free --- and forgotten.
\end{itemize}

\textbf{Adventure Hook:} A Sidhi freedman's rope has been stolen. Without it, he can be re-enslaved. The party must find the thief --- who is working for the Ashaani master who originally sold the seal. The rope has seven knots left. That is seven witnesses. That is seven days.
\end{infobox}

\chapter{The Sixth Night: The Pereshi Fire-Temple}

\section*{Told by}
A Pereshi priest of the Eternal Flame, who tended a fire that had not gone out in three hundred years. He spoke in a whisper, as if the fire were listening.

\section*{The Tale}
The Pereshi do not worship the fire. They tend it. There is a difference, they say. Worship is for gods. Tending is for guests.

A traveler --- call him Rostam --- came to a fire-temple in the mountains. He was cold, hungry, and lost. The priest let him in.

"Sit by the fire," the priest said. "But do not stare into it. The fire does not like to be watched."

Rostam sat with his back to the flame. He warmed his hands. He ate bread and salt.

"What is the fire's name?" he asked.

"It has no name," the priest said. "The fire forgets names. It remembers only warmth."

Rostam was curious. He turned and stared into the flame.

The fire was not orange. It was blue. And in the blue, he saw faces. His mother. His father. His brother, who had died in a war.

"You are looking at the dead," the priest said. "They are not angry. They are waiting. The fire is their messenger."

Rostam could not look away. He watched for an hour. When he finally turned, the priest was gone. The fire was orange again. He was alone.

He left the temple. He walked home. He never told anyone what he had seen. But he stopped being afraid of dying. He said, "The fire is warm. The dead are waiting. There is nothing to fear."

\section*{The Witch's Closing}
"Fire remembers. It does not judge. If you are not afraid of what it remembers, it will warm you. If you are, it will show you everything."

\section*{What Lurks in the Shadow of the Tale}
\begin{infobox}
\textbf{The Blue Fire} --- A sacred flame that shows visions of the dead. It cannot be used for divination --- the dead do not answer questions. They only appear. The fire is a messenger, not an oracle.

\textbf{Tags / Mechanics:} \tag{LOCATION} \tag{MEMORY} \tag{SACRED}

\begin{itemize}
\item A character who stares into the Blue Fire for an hour sees a vision of one deceased person connected to their backstory. The GM describes the vision. The vision does not answer questions --- it simply shows.
\item After seeing the vision, the character gains +1 die to resist fear for the next week (they have seen the dead and survived).
\item If a character looks away before an hour has passed, the fire goes cold. The temple must perform a ritual to relight it. The character who looked away marks 1 Corruption.
\item \textbf{Clock: The Fire's Patience [6]} -- ticks each time the Blue Fire is disturbed. When full, the fire goes out permanently, and the temple becomes a cold ruin.
\end{itemize}

\textbf{Adventure Hook:} A Pereshi fire-temple has gone dark. The priest is dead. The party must relight the Blue Fire --- but the ritual requires a memory that has never been spoken. Someone in the party must confess a secret to the ashes.
\end{infobox}

\chapter{The Seventh Night: The Weigh-Master's Grain (Oshiira)}

\section*{Told by}
A weigh-master in the Oshiiran city of Mbaro, who balanced grain for a living and balanced justice in her spare time.

\section*{The Tale}
The Oshiira say that every grain of wheat has a weight. The weight is not the same for everyone. A rich man's grain is lighter than a poor man's. The scales know. The scales are honest.

A merchant --- call him Kwame --- brought a shipment of grain to the weigh-house. The weigh-master put it on the scale. It balanced perfectly.

"The grain is light," the merchant said. "I brought more than I promised."

The weigh-master said, "The grain is not light. Your heart is heavy. You are trying to cheat."

The merchant was angry. "My heart has nothing to do with it."

The weigh-master took a single grain from the pile and placed it on the scale. The scale tipped.

"One grain of guilt," she said. "The scale weighs truth, not wheat. Bring honest grain, and it will balance. Bring lies, and it will tip."

The merchant left. He came back a week later with a new shipment. The scale balanced. He settled his promise. And he never tried to cheat again --- not because he was afraid of the law, but because he was afraid of the one grain that had tipped the scale.

\section*{The Witch's Closing}
"A scale is a witness. It does not care about your excuses. It cares about the weight of what you carry."

\section*{What Lurks in the Shadow of the Tale}
\begin{infobox}
\textbf{The Honest Scale} --- An Oshiiran weigh-house scale that measures the moral weight of goods, not their physical mass. It cannot be fooled. It will tip if the seller has cheated, even by a single grain.

\textbf{Tags / Mechanics:} \tag{ARTIFACT} \tag{TRUTH} \tag{JUDGMENT}

\begin{itemize}
\item A character who uses the Honest Scale to weigh a transaction learns whether the other party is cheating. The GM answers truthfully (yes or no). The scale does not reveal how --- only that something is wrong.
\item If the character is the one cheating, the scale's tip causes them to gain the Condition \textit{Guilty} (--1 die to all social rolls until they confess).
\item \textbf{Clock: The Unbalanced Reckoning [6]} -- ticks each time a character cheats without confessing. When full, the scale's spirit manifests and demands a public accounting. Refusal causes the character's name to be posted on the weigh-house wall as an oath-breaker.
\end{itemize}

\textbf{Adventure Hook:} A weigh-master has been murdered. The scale was found tipped, with a single grain of wheat on the guilty side. The party must discover who the grain points to --- and why they killed to hide their crime.
\end{infobox}

\chapter{The Eighth Night: The Slave Who Named the Veil (Ashaan)}

\section*{Told by}
A Breath‑less agent in the catacombs of Galanina, who spoke through a silk curtain and never showed her face.

\section*{The Tale}
In Ashaan, the Veiled Sisters bind spirits to their will. They also bind people --- but those are not called slaves. They are called "indentured." The Sisters are very careful with words.

A Sidhi man --- call him Malik --- was born into indenture. He worked the spice docks of Port Shed, loading cargo that would never be his. He learned to read and write by watching the merchants' clerks. He learned that a name is a kind of lock.

One day, a Veiled Sister came to the docks. She was looking for a scribe. Malik volunteered. He wrote letters for her. He copied contracts. He learned the names of her rivals.

The Sister was pleased. "You are useful," she said. "I will free you."

She gave him a seal --- a black disk with a single word: \textit{Mine}.

"This means you belong to no one but me," she said. "Wear it always."

Malik wore it. But the seal was warm. It whispered to him at night. It said, "You are still a slave. The word on the seal is not \textit{Mine}. It is \textit{Burden}."

Malik confronted the Sister. "You lied."

She laughed. "A seal shows what the wearer believes. You believed you were mine. So you were. Now you know the truth. Are you free?"

Malik thought. He touched the seal. It was cold now.

"I am free because I say I am," he said. "Not because of your seal."

He cut the cord and threw the seal into the sea. The Sister's face --- behind her veil --- flickered. He had seen it. Just for a moment.

"I will tell everyone what I saw," Malik said.

"No one will believe you," the Sister said.

"I do not need them to believe. I only need them to wonder."

Malik became a Breath‑less agent. He never wore a seal again. But he carried a piece of black glass in his pocket --- a shard from the seal he had thrown away. When he met a slave, he showed it to them. "This is what freedom looks like," he said. "Broken."

\section*{The Witch's Closing}
"A name is a lock. A veil is a key. Break the lock, and you are free. Lift the veil, and the Sister is afraid."

\section*{What Lurks in the Shadow of the Tale}
\begin{infobox}
\textbf{The Broken Seal} --- A shard of a Veiled Sister's seal. It no longer binds, but it can reveal the truth about any magical enslavement. When touched to a slave's collar, the collar glows red if it is bound by magic.

\textbf{Tags / Mechanics:} \tag{ARTIFACT} \tag{TRUTH} \tag{REVELATION}

\begin{itemize}
\item The Broken Seal can be used once per day to detect magical bondage on a person or object. The shard grows warm when held near a spell of binding.
\item If the shard is used to break a binding spell (by cutting the collar or speaking the slave's true name), the character who performs the ritual marks 1 Corruption (the Sister's attention turns toward them).
\item \textbf{Clock: The Sister's Gaze [8]} -- ticks each time the shard is used to free a slave. When full, a Veiled Sister appears in the character's dreams, offering a bargain: return the shard, and she will forget their name.
\end{itemize}

\textbf{Adventure Hook:} A Breath‑less agent has been captured. The party has the Broken Seal. To free her, they must enter the Sister's sanctum --- but the seal will betray their presence. Do they risk everything for one freed slave, or do they wait for a better chance?
\end{infobox}

\chapter{The Ninth Night: The Walking City's Keystone (Ngomebe)}

\section*{Told by}
An elder of the Mason‑Kin, a woman whose hands were stained with granite dust and whose voice was the sound of stone grinding on stone.

\section*{The Tale}
The Ngomebe cities walk. They move across the plateau on rollers of petrified wood, pulled by oxen and the will of the ancestors. A city that stops walking is a city that dies.

One city --- call it Khazembo --- stopped. The rollers locked. The oxen lowed. The stone‑song fell silent.

The Mason‑Kin gathered. They examined the keystone --- the heart of the city's movement. It was cracked.

"Who cracked the keystone?" the elders asked.

No one confessed. The city grew desperate. Without movement, the fields could not be rotated. The spirits grew restless.

A young mason --- call her Nia --- knelt before the keystone. "I cracked it," she said. "I hit it with my hammer. I was angry at my brother."

The elders were silent. Then the eldest spoke. "Why did you lie?"

Nia looked up. "I did not lie."

"You did," the elder said. "The keystone does not crack from a hammer. It cracks from a broken vow. You are not the breaker. You are the one who knows the breaker's name."

Nia wept. She named her brother. The brother confessed. He had sold the city's ancestral stone to a Pereshi merchant.

They retrieved the stone. The keystone healed. The city walked again.

Nia became the keystone's keeper. She never lied to the stone again. And every mason knows: a crack in the city is a crack in someone's word.

\section*{The Witch's Closing}
"A walking city is a promise on wheels. A crack is a broken vow. Find the vow, and the city will heal itself."

\section*{What Lurks in the Shadow of the Tale}
\begin{infobox}
\textbf{The Cracked Keystone} --- The heart of a Ngomebe walking city. It is a living stone that records every broken promise made within the city's bounds. It cracks when a vow is broken. It can only be healed when the oath‑breaker confesses.

\textbf{Tags / Mechanics:} \tag{LOCATION} \tag{TRUTH} \tag{JUDGMENT}

\begin{itemize}
\item A character who touches the keystone may ask, "What promise was broken here?" The keystone shows a vision of the broken vow (GM discretion). The character marks 1 Fatigue.
\item If the character is the oath‑breaker, the keystone causes them to speak the truth aloud (no roll). They cannot lie while touching the stone.
\item \textbf{Clock: The City's Death [8]} -- ticks each time a vow is broken without confession. When full, the keystone shatters, and the city becomes a ruin. The only way to reverse it is to find the original oath‑breaker's descendant and perform a ritual of apology.
\end{itemize}

\textbf{Adventure Hook:} A walking city has stopped. The keystone is cracked, but no one knows who broke it. The party must find the living descendant of the oath‑breaker --- who may not even know their ancestor's crime --- and convince them to confess.
\end{infobox}

\chapter{The Tenth Night: The Aelinnel and the Aelaerem --- The Ninth Door}

\section*{Told by}
Two voices: a gnome bell‑tender named Pim and a halfling hedge‑witch named Celandine. They argued throughout the telling, corrected each other, and eventually agreed to disagree.

\section*{The Tale}
The Aelinnel build doors that count. The Aelaerem pour cups that cannot be counted. They have lived as neighbors for centuries, borrowing each other's thresholds and arguing about the number nine.

A gnome --- call him Tink --- built a door for a halfling village. It was a fine door, made of living wood, with a bell that rang once for each opening. He counted the rings: one, two, three, four, five, six, seven, eight.

"Where is the ninth?" the halflings asked.

"There is no ninth," Tink said. "Nine is for the Hollow."

The halflings were satisfied. But a child --- call her Holly --- was curious. She rang the bell eight times, then whispered, "Nine."

The door opened onto a corridor that was not there. The corridor smelled of ash and rain. Holly walked inside.

She did not return for three days. When she did, she was carrying a stone. The stone was warm. It had a single word carved into it: \textit{Borrowed}.

The halflings blamed Tink. "Your door let the Hollow in."

Tink said, "Your child called the ninth. The door only answered."

The halflings wanted to destroy the door. Tink refused. "The door is not the problem. The problem is that we do not agree on what the ninth means. To us, nine is a number. To you, nine is a door. We need a bridge."

They built a second door, next to the first. This door had no bell. It had a cup.

"When you need to cross," Tink said, "pour water into the cup. If it fills to the brim, the ninth is safe. If it overflows, the ninth is hungry."

The halflings agreed. They kept both doors. The gnomes used the bell door. The halflings used the cup door. And every harvest, they met at the threshold and argued about whether the cup had overflowed.

\section*{The Witch's Closing}
"A bridge between cultures is not a door. It is a cup. You must agree on what full means."

\section*{What Lurks in the Shadow of the Tale}
\begin{infobox}
\textbf{The Two Doors} --- One door (gnome) rings a bell eight times and refuses the ninth. The second door (halfling) uses a cup to test the ninth. Crossing the wrong door without the proper rite invites the Hollow.

\textbf{Tags / Mechanics:} \tag{LOCATION} \tag{THRESHOLD} \tag{CULTURE}

\begin{itemize}
\item A character who uses the bell door must ring eight times. If they ring a ninth, they are lost in the Hollow for 1d3 hours and mark 1 Fatigue.
\item A character who uses the cup door must pour water. If the cup overflows, a spirit of the Hollow appears. The spirit will answer one question, but then demand a memory in return.
\item \textbf{Clock: The Threshold Dispute [6]} -- ticks each time the two cultures disagree on how to use the doors. When full, both doors slam shut. The only way to reopen them is to hold a moot where both sides tell a story that satisfies the other.
\end{itemize}

\textbf{Adventure Hook:} A gnome and a halfling village have stopped speaking. The doors are locked. The harvest is rotting. The party must arbitrate the dispute --- and find a new ritual that both cultures accept.
\end{infobox}

\chapter{The Eleventh Night: The Walking Promise}

\section*{Told by}
No one. This tale was told to me by a letter, found in the binding of this book, written in a hand I did not recognize. It is dated the same night I left the baronet's forest.

\section*{The Tale}
A promise walked into a town. It had no legs and no mouth, but it walked. It sat at the town square, and it waited.

The townspeople were afraid. "What does it want?"

An old woman said, "It wants what it is owed."

They brought coins. The promise did not take them. They brought grain. The promise did not take it. They brought a prisoner. The promise did not move.

The old woman said, "The promise is not owed in things. It is owed in stories."

She sat beside the promise. She told a story. The promise listened. When she finished, the promise stood and walked to the next town.

The townspeople were relieved. But the old woman said, "The promise will return. It always returns. A story only adds another thread."

\section*{The Witch's Closing}
"A promise is a story we tell ourselves to make the world feel fair. The world is not fair. But a story --- a story can buy time."

\section*{What Lurks in the Shadow of the Tale}
\begin{infobox}
\textbf{The Walking Promise} --- A spirit that appears when a promise has been broken and not acknowledged. It cannot be paid in coin, only in stories. It will continue to appear until the story is told.

\textbf{Tags / Mechanics:} \tag{SPIRIT} \tag{TRUTH} \tag{CURSE}

\begin{itemize}
\item The Walking Promise appears when a character breaks a significant promise and does not apologize. The GM decides when it appears --- always at a dramatic moment.
\item The Promise cannot be fought. It cannot be bargained with (except with stories). A character may tell the Promise a story. On a successful \textit{Spirit + Performance} test (DV 6), the Promise accepts the story and leaves for a week.
\item On a failure, the Promise takes something --- a memory, a skill point, or a relationship. The character may choose what is taken, but it must be significant.
\item \textbf{Clock: The Promise's Return [8]} -- ticks each week the Promise is not satisfied. When full, the Promise claims the character's name. They become a stranger to everyone who knew them.
\end{itemize}

\textbf{Adventure Hook:} A town is haunted by the Walking Promise. It appears every night at midnight, demanding a story. The townspeople are running out of stories. The party must discover what original promise the spirit is collecting --- and who broke it.
\end{infobox}

\chapter*{Epilogue: The Storyteller's Return}
\addcontentsline{toc}{chapter}{Epilogue: The Storyteller's Return}

I reached the Pereshi highlands in the spring. The snow was melting. The roads were mud. I found Kaveh's village --- a cluster of stone houses built into a cliff, with a fire-temple at the center.

Kaveh was dead. He had died the winter before, of a cough that would not stop.

His daughter met me at the gate. She had his eyes and his laugh. "You are the Tulkani," she said. "My father told me you would come. He said you carried a story for him."

I knelt. "I have carried it for years. I am sorry I am late."

She helped me up. "He is not here to hear it. But I am. Tell it to me."

I told Kaveh's tale. It was not long --- no longer than any of the tales in this book. But it was his, and I had borrowed it, and now I was returning it.

When I finished, the daughter said, "The promise is returned. But my father also said that you would leave a copy of your book in his fire-temple. He said the fire would remember the stories even if he could not."

I left the book. I do not know if the fire remembers. But I like to think it does.

\dusana{On the road back to the Dolmis, with a lighter chest. The promise is returned. The stories are still being told.}

\chapter*{Postscript: A Map of Unwritten Tales}
\addcontentsline{toc}{chapter}{Postscript: A Map of Unwritten Tales}

This is the third cycle. There will be a fourth.

I have heard whispers of a place called the Crimson Basin, where the river runs red with iron and the spirits are older than the first empire. I have heard of a desert called the Wound, where the Veiled Sisters bind slaves to their service with magic that has no name. I have heard of walking cities that move across the plateau, and of a temple where the Kalrantha amulet was first sealed.

I do not know if I will write those tales. The road is long, and the ink is expensive.

But I carry a promise to the Pereshi, and to the Sidhi, and to the Hyaar, and to all the peoples of the southern lands. I have returned some of what I borrowed. Not all.

A cycle is never closed. It is only passed to the next teller.

\dusana{At the crossroads of three empires, with a full waterskin and an empty purse. The ninth cup remains un-poured. The ninth tale remains untold.}

\chapter*{Appendix: Customs of the Southern Roads}
\addcontentsline{toc}{chapter}{Appendix: Customs of the Southern Roads}

\begin{table}[h]
\centering
\begin{tabularx}{\textwidth}{p{4cm} p{4cm} X}
\toprule
\textbf{Custom} & \textbf{Culture} & \textbf{Consequence of Breaking} \\
\midrule
Do not stare into the Blue Fire. & Pereshi (fire-temples) & You see the dead. You may not look away. \\
Turn the lamp a quarter degree to judge a ship. & Sidhi (lighthouses) & The lamp may wreck a different ship. \\
Answer the djinn's three questions truthfully. & Hyaar (high passes) & The djinn may keep your answer --- or you. \\
Tie nine knots in a rope to witness a promise. & Sidhi (freedmen) & The law hesitates. The witnesses remember. \\
Tell a story to the river before crossing. & Aeler (bridges) & The bridge may collapse. The river may take a life. \\
A carpet's border shows unsettled weights. & Pereshi (weavers) & The weights grow. The carpet unravels. \\
The weigh‑master's scale measures truth. & Oshiira (grain trade) & A single grain of guilt tips the balance. \\
A Veiled Sister's seal shows what you believe. & Ashaan (indenture) & Break it, and she will see your face. \\
A walking city's keystone cracks with every broken vow. & Ngomebe (Mason‑Kin) & The city stops. The stone must hear a confession. \\
A door that rings eight times should never answer a ninth. & Aelinnel (gnomes) & The Hollow opens. You may not return. \\
A cup that overflows invites the Hollow. & Aelaerem (halflings) & A spirit will ask for a memory. \\
The Walking Promise accepts only stories. & All & A story buys time. No story buys nothing. \\
\bottomrule
\end{tabularx}
\end{table}

These are not laws. They are courtesies extended to the spirits who share the world. Violate them, and the spirits will not punish you --- they will simply stop being polite.

And impolite spirits are the most dangerous kind.

\chapter*{Colophon}
\addcontentsline{toc}{chapter}{Colophon}

This book was written on paper made from papyrus, in ink cut with salt water and the ash of a Sidhi lighthouse lamp. The binding is rope, tied with nine knots.

The ninth knot is not cut.

The promise is returned.

The road is still south.

\dusana{At the edge of the Crimson Basin, with a new notebook and an old lamp. The ink is fresh. The tales are waiting.}

\end{document}